\subNsection{1}{Disjunktive Normalform}{Duy Tan Pham}

Die \gls{dnf} ist ein eigener Themenbereich im Modul
\glqq Grundlagen der Informatik\grqq{}. Sie folgt dem gleichen Aufbau wie die anderen
Themen: Übersicht, Übungen, Beispiele und eine freie Simulation. Die Inhalte bauen
schrittweise aufeinander auf: von der kanonischen DNF über die minimale DNF bis zur
optional global kleinsten Formel. Die einzelnen Schritte der Ableitung werden dabei
sichtbar gemacht.

Die inhaltliche Ausrichtung folgt den vorhandenen Übungsunterlagen der Lehrveranstaltung.
Dort wird die DNF aus Wahrheitstabellen mit drei und vier Variablen gebildet. Der Modus
\glqq Wahrheitstabelle nach DNF\grqq{} setzt genau diesen Aufgabentyp interaktiv um. Die
Grenze von vier Variablen deckt damit zugleich die vorgesehenen Übungsbeispiele ab.

\subNsection{2}{Entscheidungen in der Architektur}{Duy Tan Pham}

Abbildung~\ref{fig:dnf_ablauf} fasst den fachlichen Datenfluss der beiden
Eingabemodi zusammen.

\begin{figure}[htbp]
  \centering
  \begin{tikzpicture}[
    font=\footnotesize,
    box/.style={draw, rounded corners=2pt, align=center, inner sep=5pt, minimum height=1.0cm},
    flow/.style={-{Stealth[length=2.2mm,width=1.8mm]}, semithick},
    bound/.style={flow, dashed},
    lbl/.style={font=\scriptsize\itshape, inner sep=2pt, fill=white}
  ]
    \node[box, text width=2.1cm, fill=black!5]  (source)    at (0,0)
      {Boolescher\\Ausdruck};
    \node[box, text width=2.1cm, fill=black!5]  (parser)    at (3,0)
      {Parser und\\Syntaxbaum};
    \node[box, text width=2.1cm, fill=black!12] (table)     at (6,0)
      {Referenz-\\tabelle};
    \node[box, text width=2.1cm, fill=black!5]  (canonical) at (9,0)
      {Kanonische DNF};
    \node[box, text width=2.1cm, fill=black!12] (minimal)   at (12,0)
      {Minimale DNF\\exakte Minimierung};

    \node[box, text width=2.1cm, fill=black!5]  (tableinput) at (0,-2.3)
      {Vorgegebene\\Tabelle};
    \node[box, text width=2.1cm, fill=black!5]  (answer)     at (3,-4.6)
      {Antwort der\\Lernenden};
    \node[box, text width=2.5cm, fill=black!12] (check)      at (6,-4.6)
      {Prüfung von Form\\und Äquivalenz};
    \node[box, text width=2.7cm, fill=black!12] (global)     at (10.5,-3.0)
      {Optional: global\\kleinste Formel};

    \draw[flow] (source) -- (parser);
    \draw[flow] (parser) -- (table);
    \draw[flow] (table) -- (canonical);
    \draw[flow] (canonical) -- (minimal);
    \draw[flow] (tableinput.east) -- (1.5,-2.3) -| (table.south)
      node[lbl, near start, above] {direkt};
    \draw[flow] (answer) -- (check);
    \draw[flow] (table.south) -- (check.north)
      node[lbl, midway, right] {Referenz};
    \draw[flow] (table.south east) to[out=-25,in=165]
      node[lbl, midway, above] {Funktion} (global.west);
    \draw[bound] (minimal.south) to[out=-90,in=35]
      node[lbl, midway, right] {Obergrenze} (global.north east);
  \end{tikzpicture}
  \caption[Fachlicher Datenfluss im DNF-Bereich]{%
    Fachlicher Datenfluss im DNF-Bereich. Die Referenztabelle verbindet beide
    Eingabemodi mit Ableitung, Minimierung, globaler Suche und Äquivalenzprüfung.}
  \label{fig:dnf_ablauf}
\end{figure}

\begin{itemize}
  \item Der Bereich trennt \texttt{domain}, \texttt{application}, \texttt{adapters}
  und \texttt{presentation}. Parser, Wahrheitstabelle, DNF-Ableitung, Minimierung und
  Prüfung liegen in \texttt{domain} und lassen sich ohne React testen.

  \item Der boolesche Ausdruck wird von einem eigenen Parser (\texttt{dnfParser.ts}) in
  einen Syntaxbaum umgewandelt. Aus diesem Baum wird die Wahrheitstabelle erzeugt. Die
  Wahrheitstabelle ist die gemeinsame Grundlage für alle weiteren Schritte. Der Parser
  erlaubt bewusst mehrere Schreibweisen (\texttt{!}, \texttt{\&}, \texttt{|} und ebenso
  \(\lnot\), \(\land\), \(\lor\)) und meldet Fehler mit Position.

  \item Die Ableitung trennt drei Ergebnisstufen sauber: die kanonische DNF
  (\texttt{dnfDerivation.ts}, ein Vollminterm je Eins-Zeile \cite[S.~163]{schiffmann2003}), die minimale DNF
  (\texttt{quineMcCluskey.ts}) und die optional global kleinste Formel
  (\texttt{minimalExpression.ts}). Berechnung und Darstellung bleiben dabei getrennt.

  \item Die wichtigste fachliche Entscheidung betrifft die Minimierung. Die erste Version
  wählte die restliche Überdeckung nach Quine-McCluskey
  \cite[S.~172]{schiffmann2003} \emph{greedy} aus und konnte dadurch die beste Lösung
  verpassen. Quine-McCluskey bildet nun die Primimplikanten. Die Funktion
  \texttt{selectCover} wählt daraus mit Branch-and-Bound exakt eine vollständige
  Abdeckung mit minimaler Literalzahl. Zweige, die das beste bekannte Ergebnis nicht
  mehr verbessern können, werden frühzeitig verworfen.

  \item Für die optional kleinste Formel wird zusätzlich eine exakte Synthese über den
  Funktionsraum berechnet. Die minimale DNF liefert dafür eine gültige Obergrenze. Für
  Ausdrücke aus Literalen, die mit UND und ODER verbunden werden, findet die Suche bis
  vier Variablen das globale Literalminimum. Negationen stehen dabei nur direkt an
  Variablen. Da die Anzahl möglicher Funktionen mit \(2^{2^{n}}\) stark wächst, ist das
  Werkzeug auf vier Variablen begrenzt. Diese Grenze wird in Anwendungsschicht,
  Tabellen-Eingabe und Import geprüft.

  \item Die Prüfung ist von der Darstellung getrennt. Ein eingegebener Ausdruck wird für
  jede Zeile der Wahrheitstabelle ausgewertet und nicht durch einen reinen Textvergleich
  bewertet. Dadurch erkennt die Übungsprüfung (\texttt{dnfExerciseCheck.ts}) korrekte,
  nur äquivalente und nicht äquivalente Antworten. Die Vereinfachungsprüfung
  (\texttt{dnfSimplificationCheck.ts}) prüft zusätzlich, ob die Literalzahl das exakte
  Minimum erreicht.

  \item Import, Export und automatisches Speichern verwenden den gemeinsamen
  Training-Center-Speicher und einen kompakten, URL-sicheren Base64-Code.
\end{itemize}

\subNsection{2}{Entscheidungen im UI}{Duy Tan Pham}

\begin{itemize}
  \item Der DNF-Bereich bietet zwei Aufgabenmodi: \glqq Ausdruck nach
  Wahrheitstabelle\grqq{} und den umgekehrten Weg \glqq Wahrheitstabelle nach DNF\grqq{}.
  Beim zweiten Modus stellen die Lernenden aus einer gegebenen Tabelle die DNF selbst auf.
  Beide Übungen nutzen dieselben Bausteine.

  \item Die einzelnen Schritte der Ableitung werden sichtbar gemacht: von den Eins-Zeilen
  der Tabelle über die passenden kanonischen Terme bis zur minimalen Form. Bei der
  Vereinfachung wird nicht nur das Ergebnis gezeigt, sondern auch der Lösungsweg. In den
  Übungen bearbeiten die Lernenden Wahrheitstabellen oder DNF-Ausdrücke selbst und
  erhalten eine fachlich berechnete Rückmeldung.

  \item Kleinere Wahrheitstabellen werden vollständig angezeigt. Bei \texttt{16} Zeilen,
  also bei vier Variablen, gibt es eine kompakte, scrollbare Ansicht, die sich auf eine
  Vollansicht umschalten lässt. Dadurch bleibt die maximale Darstellungshöhe begrenzt.

  \item Folgeaktionen bleiben sichtbar und werden bis zum Erfüllen ihrer Voraussetzungen
  ausgegraut. Nur kontextabhängige Schaltflächen, etwa für die Vollansicht einer großen
  Tabelle, erscheinen bei Bedarf. Dadurch bleibt die Bedienung vorhersehbar.

  \item Eingaben werden vereinheitlicht: Lernende dürfen \texttt{!}, \texttt{\&} und
  \texttt{|} tippen, angezeigt wird aber immer \(\lnot\), \(\land\), \(\lor\). Das macht die
  Eingabe leichter, ohne die Darstellung uneinheitlich zu machen.

  \item Der DNF-Bereich nutzt den gemeinsamen Themenaufbau mit Kopfbereich, Tabs und
  globaler Sidebar. Dadurch fügt er sich in die übrige Lernplattform ein.
\end{itemize}

\subNsection{2}{Andere Entscheidungen}{Duy Tan Pham}

Die kanonische und die minimale DNF werden bewusst getrennt dargestellt. In der
kanonischen DNF enthält jeder Minterm alle Variablen. Eine minimale DNF kann dagegen
kürzere Terme besitzen, etwa \((a \land c) \lor (b \land c)\). Die getrennte Darstellung
macht diesen fachlichen Unterschied nachvollziehbar und verhindert, dass eine korrekte
Vereinfachung fälschlich als unvollständig angesehen wird.
